% Options for packages loaded elsewhere
\PassOptionsToPackage{unicode}{hyperref}
\PassOptionsToPackage{hyphens}{url}
%
\documentclass[
  11pt,
]{article}
\usepackage{amsmath,amssymb}
\usepackage{iftex}
\ifPDFTeX
  \usepackage[T1]{fontenc}
  \usepackage[utf8]{inputenc}
  \usepackage{textcomp} % provide euro and other symbols
\else % if luatex or xetex
  \usepackage{unicode-math} % this also loads fontspec
  \defaultfontfeatures{Scale=MatchLowercase}
  \defaultfontfeatures[\rmfamily]{Ligatures=TeX,Scale=1}
\fi
\usepackage{lmodern}
\ifPDFTeX\else
  % xetex/luatex font selection
\fi
% Use upquote if available, for straight quotes in verbatim environments
\IfFileExists{upquote.sty}{\usepackage{upquote}}{}
\IfFileExists{microtype.sty}{% use microtype if available
  \usepackage[]{microtype}
  \UseMicrotypeSet[protrusion]{basicmath} % disable protrusion for tt fonts
}{}
\makeatletter
\@ifundefined{KOMAClassName}{% if non-KOMA class
  \IfFileExists{parskip.sty}{%
    \usepackage{parskip}
  }{% else
    \setlength{\parindent}{0pt}
    \setlength{\parskip}{6pt plus 2pt minus 1pt}}
}{% if KOMA class
  \KOMAoptions{parskip=half}}
\makeatother
\usepackage{xcolor}
\usepackage[margin=1in]{geometry}
\setlength{\emergencystretch}{3em} % prevent overfull lines
\providecommand{\tightlist}{%
  \setlength{\itemsep}{0pt}\setlength{\parskip}{0pt}}
\setcounter{secnumdepth}{-\maxdimen} % remove section numbering
\ifLuaTeX
  \usepackage{selnolig}  % disable illegal ligatures
\fi
\usepackage{bookmark}
\IfFileExists{xurl.sty}{\usepackage{xurl}}{} % add URL line breaks if available
\urlstyle{same}
\hypersetup{
  pdftitle={P02.1: residue cardinality at a relative boundary and analytic logarithmic series},
  pdfauthor={Research draft; no independent correctness verdict},
  hidelinks,
  pdfcreator={LaTeX via pandoc}}

\title{P02.1: residue cardinality at a relative boundary and analytic
logarithmic series}
\author{Research draft; no independent correctness verdict}
\date{8 October 2026}

\begin{document}
\maketitle

\subsection{Status and exact scope}\label{status-and-exact-scope}

The original P02.1 characterization is unresolved in this manuscript. We
prove two precisely delimited results, with full proofs below. First,
for one actual classical psh boundary-jump function on a strict analytic
polydisc, local sequential semipositive real-PL regularization exists if
and only if the residue field is countable. The negative direction
excludes even nonmonotone pointwise continuous sequences. A decreasing
semipositive real-PL net always exists. This extends the supplied
Laurent-field example without using its proof as a premise: there is no
coefficient-field, algebraic-closedness, discrete-valuation, or
prescribed valuation-normalization assumption. Second, we construct one
decreasing semipositive real-PL sequence for every negative analytic
logarithmic series with positive real coefficients on a good strict
space. These series can have infinitely many noninvariant polar divisors
and value \(-\infty\). Their input description is analytic and does not
refer to an approximation class.

Neither result treats every classical psh germ. Countability is sharp
only for the specified boundary jump, and the logarithmic-series
hypothesis is an extra starting-function hypothesis. In particular, we
do not characterize all spaces having the original property. We make no
priority claim for these elementary extensions. No earlier research note
or book theorem is a mathematical premise.

\subsection{Definitions and authenticated published
inputs}\label{definitions-and-authenticated-published-inputs}

Throughout \(K\) is a complete nontrivially valued non-Archimedean
field, \(K^\circ\) its valuation ring, \(K^{\circ\circ}\) its maximal
ideal, and \(k=K^\circ/K^{\circ\circ}\) its residue field. All
neighborhoods are ordinary open neighborhoods. Fix \(a\in K\) with
\(0<|a|<1\) and put \(c=-\log|a|>0\). We use \(\log 0=-\infty\).

Classical psh means the upper semicontinuous class of {[}GR-C,
Definition 4.1{]}: after every analytic field extension \(L/K\) and
every analytic morphism from an \(L\)-analytic curve to \(X_L\), the
pullback is subharmonic in {[}GR-C, Definition 3.4{]}. The latter tests
smooth connected strict curves over nontrivially valued analytic
extensions, allowing the identically \(-\infty\) function. An analytic
curve there is separated and pure one dimensional. Thus every scalar
extension in this definition is retained; current positivity is not
substituted for the definition.

A real-PL function is continuous and locally a finite real linear
combination of integral PL functions, on the G-topology generated by
\textbf{strict} analytic domains. Integral PL functions are G-locally
logarithms of analytic units {[}GJR, Definition 5.1 and Remarks 5.2{]}.
For a real-PL function \(h\) and a point \(z\), its residue line bundle
\(L_h(z)\) lies in
\(\operatorname{Pic}(\operatorname{red}(X,z))\otimes\mathbb R\).
Semipositivity means this class is nef. Here nefness is geometric: on a
finite-type model of the reduction of the germ, every integral proper
closed residue curve has nonnegative degree. {[}GJR, Proposition 7.1 and
Definition 7.2{]} identify this condition with the existence, near the
point, of a formal-model representative whose special-fibre line bundle
is nef. It is not a definition in terms of limits or curve-test
approximation.

We apply the following published statements with their indicated scopes.

\begin{enumerate}
\def\labelenumi{\arabic{enumi}.}
\tightlist
\item
  {[}GR-C, Proposition 4.2(2),(3){]} says that logarithms of analytic
  absolute values are classical psh, and that nonnegative real linear
  combinations and finite maxima preserve classical psh. The extended
  convention is included. {[}GR-C, Proposition 4.2(4){]} gives
  decreasing-net closure. We also supply a direct harmonic-comparison
  argument for the finite approximants used here below.
\item
  {[}GR-C, Theorem 4.3{]} says: over a nontrivially valued complete
  field, on a good strictly analytic space, an \textbf{already
  real-valued real-PL} function is semipositive if and only if it is
  classical psh. Its statement has no extra separatedness,
  pure-dimensionality, or boundaryless assumption. We use it only after
  separately proving strict PL membership. It does not assert
  regularization of arbitrary classical psh functions.
\item
  {[}GR-C, Section 2.4{]} gives the evaluation topology of
  \(\mathcal M(A)\) and its compact Hausdorff property. {[}GR-C, Section
  3.1(3){]} identifies the relative boundary of an affinoid curve with
  its Shilov boundary. The boundary composition rule {[}GR-C, Section
  2.11(4'){]} is
  \[\partial(X/S)\subseteq\partial(X/Y)\cup f^{-1}\partial(Y/S),\] with
  equality if \(Y\to S\) is locally separated. We use this equality with
  the separated closed unit disc as \(Y\).
\item
  {[}GJR, Lemma 5.4(1),(2){]} preserves real-PL membership under
  morphism and analytic scalar pullback between good strict spaces.
  {[}GJR, Proposition 7.6(1),(2){]} preserves semipositivity in the same
  directions. No converse or normalization descent is used here.
\end{enumerate}

The primary sources are distinct. {[}GR-C{]} is \emph{Classically psh
and pluriharmonic functions on Berkovich spaces}, arXiv:2506.13548v2, 17
September 2025. {[}GR-PL{]} is \emph{On semipositive piecewise linear
functions in non-archimedean analytic geometry}, arXiv:2503.05513v1, 7
March 2025. {[}GJR{]} is \emph{Forms on Berkovich spaces based on
harmonic tropicalizations}, arXiv:2111.05741v3, 7 March 2025. Their
title/version headers in retained primary PDFs agree with the public
arXiv records.

Here is the proof dependency of the comparison, rather than a citation
to a book label. The proof of {[}GR-C, Theorem 4.3{]} first uses the
curve comparison in {[}GR-PL, Proposition 5.1{]} and semipositivity
pullback in {[}GJR, Proposition 7.6{]}. Its converse uses smooth-curve
recognition in {[}GR-PL, Corollary 5.4{]}. The June paper calls those
March results Proposition 3.1 and Corollary 3.5; the retained March
version actually numbers them 5.1 and 5.4. Proposition 5.1 compares the
residue degree and the curve Laplacian, then passes to the completed
algebraic closure. Corollary 5.4 uses Theorem 5.3 to lift proper residue
curves to analytic curves, checks them after algebraic closure and
normalization, and descends semipositivity through proper surjective
morphisms. Boundary points of the test curves impose no
proper-residue-curve condition. These arguments explain why the applied
theorem includes the boundary points needed below.

The complete immediate statements and proofs of these applied comparison
and pullback results were read in retained primary text. We import their
published foundations: Thuillier's harmonic comparison and
Poincare--Lelong results, Wanner's PL/current comparison as used in
{[}GR-PL, Proposition 5.1{]}, the local lifting theorem {[}GR-PL,
Theorem 4.1{]}, formal-model existence, Temkin reductions, and degree
projection formulas. We do not claim a fresh recursive proof of every
one of those foundations. The application ledger states this boundary
explicitly. In particular, no ordinary CLD-smooth conclusion, no locally
Fubini--Study identification of arbitrary real-PL functions, and no
removal of Wanner's smooth-regularization hypotheses is inferred.

\subsection{Strict integral truncations and decreasing
limits}\label{strict-integral-truncations-and-decreasing-limits}

\textbf{Lemma 1 (legal truncation).} Let \(X\) be good and strictly
\(K\)-analytic, let \(f\) be analytic on \(X\), and let
\(b\in K^\times\). The finite function \[H_{b,f}=\log\max\{|b|,|f|\}\]
is integral PL and semipositive on \(X\). Positive real multiples and
finite positive real sums of such functions are semipositive real-PL.

\textbf{Proof.} On each strict affinoid neighborhood, the two strict
Laurent domains \(|f|\le |b|\) and \(|f|\ge |b|\) form a finite G-cover.
On the first the expression is \(\log|b|\), the logarithm of a constant
unit. On the second \(f\) is a unit: its inverse is supplied by the
Laurent-domain algebra adjoining \(b/f\). There the expression is
\(\log|f|\). On the intersection both expressions have value
\(\log|b|\). Consequently these are the local expressions of one
continuous global function, and prove strict integral PL membership,
including at zeros of \(f\). The domains are strict because the cut is
at an actual \(K\)-value, not at an arbitrary real radius.

The constant \(\log|b|\) is classical psh, as is \(\log|f|\) with its
permitted poles. Their maximum is classical psh by the log/max
properties stated above. The real-PL comparison now proves residue-nef
semipositivity on all of \(X\). For positive real sums, refine the
finitely many strict covers and add their unit-log expressions. The sum
is real-PL. The classical cone property and the same comparison prove
semipositivity. This does not use closure under maxima of arbitrary
real-PL functions. Such closure can fail in the strict topology, as
{[}GJR, Example 5.3{]} demonstrates. \(\square\)

\textbf{Lemma 2 (decreasing limits in the present constructions).} The
infimum of a decreasing net of finite classical psh functions is
classical psh, provided its values are in \(\mathbb R\cup\{-\infty\}\).

\textbf{Proof.} The infimum is upper semicontinuous since its strict
sublevel sets are unions of the strict sublevel sets of the
approximants. Pull the whole net back after any field extension and then
to a smooth connected strict test curve \(C\). Such smooth tests suffice
by {[}GR-C, Section 4.1.3{]}. Denote the finite subharmonic pullbacks by
\(g_i\) and their infimum by \(g\). If \(g\equiv-\infty\) it is allowed.
Otherwise use Thuillier's harmonic-comparison definition recalled in
{[}GR-C, Section 3.3{]}. Let \(A\subset C\) be strictly affinoid and
\(H\) harmonic on \(A\), with \(g\le H\) on its finite boundary. For
\(\epsilon>0\), at each boundary point choose an index with
\(g_i<H+\epsilon\). A single upper bound for the finitely many chosen
indices makes this inequality hold on the entire boundary. Harmonic
comparison for that \(g_i\) gives \(g_i\le H+\epsilon\) on \(A\). Taking
the infimum and then \(\epsilon\downarrow0\) gives \(g\le H\) on \(A\).
Thus \(g\) is subharmonic, including its extended negative values. Every
allowed smooth test has been checked. The general statement agrees with
{[}GR-C, Proposition 4.2(4){]}. \(\square\)

\subsection{The Gauss boundary jump over an arbitrary
field}\label{the-gauss-boundary-jump-over-an-arbitrary-field}

Put \[D=\mathcal M(K\langle T\rangle),\qquad
X_n=\mathcal M(K\langle T_1,\ldots,T_n\rangle),\quad n\ge2.\] Let
\(\eta\) be the Gauss point of \(D\), \(p:X_n\to D\) the
first-coordinate projection, and \(i:D\to X_n\) the analytic section
given by \(T_1=T\) and \(T_2=\cdots=T_n=0\). Define
\[w(y)=\begin{cases}1&y=\eta,\\0&y\ne\eta,\end{cases}
\qquad u=w\circ p,\qquad x=i(\eta).\] These are actual strict affinoid
spaces. They are good and separated, and the standard Tate-algebra
dimension theorem gives dimensions \(1\) and \(n\). In particular the
example is genuinely higher dimensional, not an abstract tropical fan.
Relative boundary is retained.

\textbf{Lemma 3 (reduction and detection).} For each monic irreducible
polynomial \(q\in k[T]\), choose a monic lift \(Q_q\in K^\circ[T]\).
Write \(\mathcal R\) for the set of all such \(q\). For every
\(y\ne\eta\) in \(D\) there exists \(q\in\mathcal R\) with
\(|Q_q(y)|<1\).

\textbf{Proof.} Each nonzero restricted series \(f=\sum b_jT^j\) has
attained Gauss norm \(\|f\|=\max_j|b_j|\in|K^\times|\), because
\(b_j\to0\). After dividing by an attaining coefficient, its reduction
is a nonzero polynomial in \(k[T]\). Reduction of products proves that
the Gauss norm is multiplicative, since \(k[T]\) is a domain. It defines
\(\eta\).

Every bounded multiplicative seminorm on the Tate algebra is bounded by
its Gauss norm. For completeness, if boundedness gives
\(|h(y)|\le C\|h\|\), apply it to \(h^N\) and take \(N\)th roots,
letting \(N\to\infty\). If \(y\ne\eta\), some \(f\) has
\(|f(y)|<\|f\|\). Normalize to \(\|f\|=1\). Factor its reduction as
\[\bar f=\bar d\prod_{r=1}^s q_r^{e_r},\qquad \bar d\in k^\times.\]
Choose a unit lift \(d\) of \(\bar d\) and set
\(R=d\prod Q_{q_r}^{e_r}\). Then \(\|f-R\|<1\). This strict inequality
also holds for nondiscrete valuations: the coefficients of the
difference have norms less than one and tend to zero, so their maximum
is attained and less than one. Hence \(|R(y)|<1\). Each factor
\(|Q_{q_r}(y)|\) is at most one and \(|d|=1\), so some factor is less
than one. A nonzero constant reduction would instead force \(|f(y)|=1\)
and is excluded. No algebraic closedness or coefficient-field lifting is
used. \(\square\)

\textbf{Theorem 4 (exact sequential criterion for this jump).} For every
such \(K\) and every \(n\ge2\), \(u\) is finite bounded upper
semicontinuous and classical psh, and \(x\in\partial_KX_n\). The
following are equivalent:

\begin{itemize}
\tightlist
\item
  \(k\) is countable;
\item
  \(u\) has a decreasing sequence of finite semipositive real-PL
  approximants on all of \(X_n\);
\item
  some ordinary neighborhood of \(x\) admits a sequence of continuous
  real functions converging pointwise to \(u\).
\end{itemize}

Regardless of the cardinality of \(k\), \(u\) has a decreasing
semipositive real-PL net on all of \(X_n\).

\textbf{Proof: the net and classical admissibility.} Direct
\[I=\operatorname{Fin}(\mathcal R)\times\mathbb N_{\ge1}\] by inclusion
in the first coordinate and the usual order in the second. For
\((F,m)\in I\) put \[P_F=\prod_{q\in F}Q_q,\quad P_\varnothing=1,\qquad
V_{F,m}=\frac1c\log\max\{1,|a^{-1}P_F^m|\}
=\max\{0,1+(m/c)\log|P_F|\}.\] The first formula, including at zeros,
proves finite semipositive real-PL membership by Lemma 1 with \(b=1\)
and the analytic function \(a^{-1}P_F^m\), followed by positive real
scaling. The second formula is a numerical identity, not an assertion of
arbitrary-real-max closure.

Since every \(Q_q\) has norm at most one on \(D\), increasing \(F\)
decreases \(|P_F|\), and increasing \(m\) decreases the second formula.
At \(\eta\), every \(P_F\) is monic with Gauss norm one, so
\(V_{F,m}(\eta)=1\). For any other point, Lemma 3 provides a detecting
factor \(Q_q\) with norm \(r<1\). At every index containing that factor,
\(V_{F,m}\le\max\{0,1+(m/c)\log r\}\), which is zero for sufficiently
large \(m\); if \(r=0\) it is zero immediately. Therefore
\(V_{F,m}\downarrow w\) pointwise. Replace \(T\) by \(T_1\) to obtain
exactly the same net on \(X_n\), decreasing to \(u\). Lemma 2 proves
classical psh of both targets with all scalar curve tests. This proof
does not infer the classical class from current positivity or from an
old boundary audit.

The singleton \(\{\eta\}\) is closed in \(D\), so \(w\) is upper
semicontinuous, and consequently so is \(u\). The Gauss point attains
the norm of every analytic function, and hence its singleton is the
Shilov boundary. The relative boundary of the affinoid curve \(D\) is
therefore \(\{\eta\}\). Apply the boundary composition equality to
\(X_n\xrightarrow p D\to\mathcal M(K)\); the structural map of \(D\) is
separated. It gives \(p^{-1}(\eta)\subseteq\partial_KX_n\), in
particular \(x\in\partial_KX_n\).

\textbf{Proof: the countable direction.} If \(k\) is countable, its
polynomial ring is countable, so \(\mathcal R\) is countable. It is
infinite: if only finitely many monic irreducibles existed, their
product plus one would have a new irreducible factor. Enumerate
\(\mathcal R=\{q_1,q_2,\ldots\}\) and take \(F_j=\{q_1,\ldots,q_j\}\),
\(m_j=j\). This sequence is cofinal in \(I\): for any finite \(F\) and
integer \(m\), all sufficiently large \(j\) contain \(F\) and satisfy
\(j\ge m\). It is therefore decreasing with the same infimum. This is an
explicit cofinality proof, not an extraction from metrizability. A
countable dense subset of \(K\) is unnecessary for this particular
function. Its pullback sequence on \(X_n\) proves the second assertion,
which implies the third by continuity.

\textbf{Proof: the uncountable obstruction.} Choose one representative
\(b_\alpha\in K^\circ\) for every residue class \(\alpha\in k\), and let
\(z_\alpha\in D\) be the rigid point \(T=b_\alpha\). Any ordinary
neighborhood of \(\eta\) contains all but finitely many \(z_\alpha\).
Indeed, for a nonzero \(f\in K\langle T\rangle\), normalize by an
attaining coefficient. The normalized reduction is a nonzero polynomial.
Except at its finitely many roots in \(k\), evaluation at \(b_\alpha\)
is a unit, so \(|f(z_\alpha)|=|f(\eta)|\). A basic neighborhood imposes
conditions on finitely many analytic evaluations, and the union of the
exceptional root sets is finite. The assertion follows from the
evaluation topology.

Suppose \(h_j\) are continuous real functions on an ordinary
neighborhood \(W\) of \(x\) and converge pointwise to \(u|_W\). Let
\(U=i^{-1}(W)\), a neighborhood of \(\eta\). For each \(j\), continuity
gives a neighborhood \(N_j\subset U\) of \(\eta\) such that
\[|h_j(i(y))-h_j(x)|<1/4\quad(y\in N_j).\] Each \(N_j\) misses only
finitely many \(z_\alpha\). Uncountability provides one \(\alpha\)
outside the union of all these exceptional finite sets. Then
\(i(z_\alpha)\in W\) and the displayed inequality holds for every \(j\).
But \(h_j(x)\to1\) and \(h_j(i(z_\alpha))\to0\), a contradiction. The
two points are distinct because \(T-b_\alpha\) has norm zero at
\(z_\alpha\) and norm one at \(\eta\). Monotonicity was never used. Thus
the third assertion fails when \(k\) is uncountable, finishing the
equivalence. \(\square\)

The same negative argument holds on any good strict space \(Y\) with
analytic maps \(D\xrightarrow{s}Y\xrightarrow{r}D\),
\(r\circ s=\mathrm{id}\), for the function \(w\circ r\) at \(s(\eta)\).
It is classical psh by pullback. Pulling a hypothetical continuous
approximating sequence back through \(s\) contradicts the proof above.
The polydiscs give actual examples of this geometric retraction
criterion; it is not a classification of all spaces.

\subsection{Scalar extension: the native jump and the pulled-back jump
differ}\label{scalar-extension-the-native-jump-and-the-pulled-back-jump-differ}

For any analytic extension \(L/K\), the same formulas \(V_{F,m}\), with
the same \(K\)-defined polynomials, on \(D_L\) converge to
\(w\circ\pi\), where \(\pi:D_L\to D\). The reason is literal equality of
the norms of a \(K\)-defined function at a point upstairs and its image
downstairs. Their semipositivity and real-PL membership either follow
from {[}GJR, 5.4 and 7.6{]} or from the same strict Laurent-cover proof
over \(L\). The corresponding statement on \((X_n)_L\) is identical. All
the functions are defined on one whole domain, and their formulas agree
on every domain intersection. No independent chartwise choices are
glued.

The support of \(w\circ\pi\) is \(\pi^{-1}(\eta)\), which need not be a
singleton. It must not be identified with the native jump
\(w_L=\mathbf1_{\{\eta_L\}}\) on \(D_L\). Theorem 4 applied over \(L\)
says that the native jump has a continuous sequential approximation near
its Gauss point if and only if \(\widetilde L\) is countable. This is
compatible with pullback of a sequence for \(w\).

Here is an explicit instance. Set \(K=\mathbb Q_p\). Let \(A\) be an
uncountable set, put \(F=\mathbb Q_p(t_\alpha:\alpha\in A)\), and give
each polynomial the Gauss \(p\)-adic valuation (the minimum coefficient
valuation); extend it to fractions. Polynomial reduction in the domain
\(\mathbb F_p[t_\alpha]\) proves multiplicativity, so this is a valued
field extending \(\mathbb Q_p\). Each polynomial uses only finitely many
variables. Its residue field is \(\mathbb F_p(t_\alpha:\alpha\in A)\):
normalize numerator and denominator to coefficient valuation zero to
compute residues, and lift any residue rational function. Let \(L\) be
the completion. Completion does not change this residue field, since any
element of norm at most one is within distance less than one of an
element of \(F\) and thus has the same residue. The residue field of
\(L\) is uncountable, while that of \(K\) is finite.

Thus \(w_K\) has a decreasing real-PL sequence, while \(w_L\) has no
continuous sequential approximation near \(\eta_L\). For a chosen
\(\alpha\), the \(L\)-point \(T=t_\alpha\) maps to \(\eta_K\). To see
this, normalize any nonzero \(K\)-Tate series to norm one; its reduction
is a nonzero polynomial in \(\mathbb F_p[T]\), which remains nonzero
when evaluated at the transcendental residue of \(t_\alpha\). Hence its
value at that point is its Gauss norm. Consequently
\((w_K\circ\pi)(t_\alpha)=1\), whereas \(w_L(t_\alpha)=0\). No assertion
that the full P02.1 property holds over \(\mathbb Q_p\) is made.

\subsection{An explicit sector with extended, noninvariant
poles}\label{an-explicit-sector-with-extended-noninvariant-poles}

\textbf{Theorem 5 (analytic logarithmic series).} Let \(X\) be any good
strictly \(K\)-analytic space. Let \(f_1,f_2,\ldots\) be analytic
functions on this one space with \(|f_i|\le1\) everywhere, and let
\(\lambda_i>0\) be any real numbers. Define the nonpositive extended sum
\[u(y)=\sum_{i=1}^\infty\lambda_i\log|f_i(y)|
:=\inf_N\sum_{i=1}^N\lambda_i\log|f_i(y)|.\] No summability or
finiteness is assumed. The function is classical psh, and the finite
semipositive real-PL functions
\[u_j=\sum_{i=1}^j\lambda_i\log\max\{|a|^j,|f_i|\}\] decrease pointwise
to \(u\) on all of \(X\). The same sequence pulls back under every
analytic field extension and every morphism from a good strict space.
Where \(u\) is finite and continuous, convergence is uniform on compact
subsets.

\textbf{Proof.} Each summand of \(u_j\) is finite and nonpositive. Lemma
1, with \(b=a^j\), proves semipositive real-PL membership. Finite common
refinements of its strict Laurent covers prove membership for the sum.
These refinements require no maxima of real-PL functions or cuts at
irrational values.

Increasing \(j\) decreases the truncation threshold in each of the first
\(j\) summands and adds one nonpositive summand. Thus \(u_{j+1}\le u_j\)
pointwise. Omitting a nonpositive tail and raising individual summands
by truncation gives \(u_j\ge u\). If \(\ell=\inf_j u_j\), then for each
fixed \(N\) and every \(j\ge N\),
\[u_j\le\sum_{i=1}^N\lambda_i\log\max\{|a|^j,|f_i|\}.\] Let
\(j\to\infty\) at a fixed point. The finite sum on the right decreases
to \(\sum_{i=1}^N\lambda_i\log|f_i|\), including value \(-\infty\) if a
factor vanishes. Hence \(\ell\le\sum_{i=1}^N\lambda_i\log|f_i|\) for
every \(N\), and therefore \(\ell\le u\). Together with \(\ell\ge u\),
this proves equality at every point, also when \(u=-\infty\). Lemma 2
proves classical psh of the limit. Zero coefficients, if desired, are
simply discarded; we never form \(0\cdot(-\infty)\).

Under a morphism or scalar projection, every displayed evaluation and
every truncation has the same value as the pullback. Thus monotonicity
and convergence are unchanged pointwise. The pulled-back analytic
functions remain bounded by one, and the same strict-cover construction
and comparison apply on a good strict source. The statement concerns
pullbacks of these input data, not arbitrary new functions over an
extension field.

On a compact subset \(C\) where \(u\) is finite continuous, for any
\(\epsilon>0\) the open subsets of \(C\) given by \(u_j-u<\epsilon\)
increase and cover \(C\). Compactness supplies a single index covering
\(C\), proving uniform convergence. This argument needs no
metrizability. For discontinuous or extended-valued limits we assert
only pointwise convergence. \(\square\)

For a concrete noninvariant example take \(K=\mathbb Q_p\) and the
strict bidisc. Set \(f_1=T_1-T_2\), \(f_i=T_1-p^{i-1}T_2\) for
\(i\ge2\), and \(\lambda_i=2^{-i}\). The Gauss point has \(u=0\). At the
rigid point \((1,1)\) the first factor vanishes, so \(u=-\infty\). At
\((1+p,1)\) the first factor has norm \(|p|\) and every later factor has
norm one, so \(u=(1/2)\log|p|\). Both coordinates have norm one at both
points. The unit torus element \((1+p,1)\) takes the first point to the
second and changes \(u\), so this function is not invariant under the
maximal compact torus. All the \(f_i\) have norm at most one on the
entire bidisc. Theorem 5 supplies an actual sequence on that whole
domain, including all its polar lines and relative boundary points.

\subsection{What these proofs do and do not
establish}\label{what-these-proofs-do-and-do-not-establish}

Theorem 4 is a sharp criterion for the specified native boundary-jump
function. It excludes universal \textbf{sequential} regularization on
good strict spaces over uncountable residue fields when these boundary
points are allowed. It does not exclude regularization by arbitrary
decreasing nets: its explicit net proves the opposite for this same
example. Countable residue field is not asserted sufficient for every
classical psh function or every space. No boundaryless counterexample or
characterization follows.

Theorem 5 proves a starting-function sector given by analytic functions
and real coefficients. It includes extended poles, noninvariant analytic
branches, and spaces with boundaries; it imposes no residue-nefness
condition on \(u\) as a definition. Nevertheless the series
representation is an extra hypothesis. There is no proof here that
arbitrary classical psh functions admit such a representation. The
theorem is therefore not an answer to the word \textbf{every} in P02.1.

The retired continuous invariant toric construction neither supplies
these input classes nor addresses this discontinuous jump. The supplied
projective-bundle argument requires domains containing entire proper
fibres, and does not extend to arbitrary fibre charts. We do not use
either historical result. Likewise, curvewise regularization does not
produce a compatible higher-dimensional neighborhood without further
work. The smooth-curve weakly smooth theorem {[}GJR, Proposition
15.11{]} supplies a net on relatively compact opens of a smooth
boundaryless curve; it is not a theorem of strict real-PL sequential
regularization. Its proof allows arbitrary skeletal subdivision points.
Ordinary CLD smoothness and weak smoothness remain distinct. Our
positive conclusions are exclusively semipositive real-PL; no additional
atom is counted.

The recent paper {[}Wu{]} defines its local psh class in Definition 3.8
by decreasing sequences of tropically convex PL functions on opens in
projective analytifications. This definition, even with its subsequent
comparisons to other generated psh classes, does not supply a theorem
identifying \textbf{every} extended classical curve-test psh function
with that class on every good strict space. We do not use any of its
theorems. Its title and version were authenticated and its immediate
local definition was read; a whole-paper correctness verdict is not
claimed.

The remaining tasks include a noncircular geometric classification,
arbitrary higher-dimensional noninvariant germs without analytic-series
hypotheses, boundaryless success/failure criteria, general
singular-space formal-model compatibility, and a countable-extraction
argument wherever a sequence rather than a net is intended. We provide
no descent of arbitrary approximants through normalization, no global
gluing of independently chosen local sequences, and no trivial-valuation
theorem. The displayed formulas solve their overlap and scalar issues on
one actual domain; they do not solve those issues for arbitrary inputs.
Accordingly \textbf{candidate=false}.

\subsection{Reproducibility and evidence
limits}\label{reproducibility-and-evidence-limits}

The standalone \texttt{P02.1-paper.tex} is generated from this
full-proof companion and compiled locally.
\texttt{programs/reproduce.py} records manifest hashes, retained
primary-source hashes, the environment, exact commands, finite
arithmetic diagnostics, and build results. Run
\texttt{python\ programs/reproduce.py} from this directory. Full primary
PDFs/texts are stored under \texttt{LOCALONLY/}, excluded from
publication by \texttt{.gitignore}; the original bound sources are never
executed. The finite checks exercise the monotonicity formulas on
explicitly bounded rational inputs. They cannot test uncountability, all
analytic curves, residue nefness, or the universal original target.
Compilation is an artifact check, not a mathematical verdict. No
supervised compute facility is needed for this bounded work, and none
was used. The supplied historical PASS receipts are data only. No
independent reviewer was invoked for this draft.

\subsection{References}\label{references}

{[}GR-C{]} Walter Gubler and Joseph Rabinoff, \emph{Classically psh and
pluriharmonic functions on Berkovich spaces},
\href{https://arxiv.org/abs/2506.13548v2}{arXiv:2506.13548v2} (17
September 2025).

{[}GR-PL{]} Walter Gubler and Joseph Rabinoff, \emph{On semipositive
piecewise linear functions in non-archimedean analytic geometry},
\href{https://arxiv.org/abs/2503.05513v1}{arXiv:2503.05513v1} (7 March
2025).

{[}GJR{]} Walter Gubler, Philipp Jell, and Joseph Rabinoff, \emph{Forms
on Berkovich spaces based on harmonic tropicalizations},
\href{https://arxiv.org/abs/2111.05741v3}{arXiv:2111.05741v3} (7 March
2025).

{[}Wu{]} Lyuhui Wu, \emph{On the local theory of non-archimedean psh
functions},
\href{https://arxiv.org/abs/2609.27610v1}{arXiv:2609.27610v1} (23
September 2026). Used for definition/scope comparison only.

{[}Ber90{]} Vladimir G. Berkovich, \emph{Spectral theory and analytic
geometry over non-Archimedean fields}, Mathematical Surveys and
Monographs 33, AMS, 1990. Standard affinoid and Tate-algebra
foundations; deeper dependencies are disclosed above.

\end{document}
